\section{冲突规范和准据法}
\begin{enumerate}
  \item \qd{学国际私法的基础的两个概念}。
  \item 国际私法的调整对象。涉外民商事关系，主体 \hz 客体 \hz 内容\zd{任何之一涉外}。
  \item 调整方式。直接调整，间接调整（冲突法调整）。
  \item 法律适用规范就是\zd{冲突规范}（重点的重点）。
  \item 自然人的\zd{经常居所地}。作为生活中心\zd{连续居住 1 年以上}（\al{住酒店}），\zd{就医、劳务、派遣公务}等除外。
  \item 法人的国籍 = \zd{登记设立地}，住所 = \zd{主要办事机构所在地} \yjt 主营业地 \yjt（多个）最密切联系。
  \item 属人法。当事人的本国法和住所地法。
  \item 本国法。当事人国籍所属国的法律。
  \item 仅指实体法，不包括冲突法（我国没有反致），程序法（适用法院地法）
\end{enumerate}

\begin{kk}{冲突规范（法律适用规范）}
  \begin{enumerate}
    \item 和准据法（实体法）。导航和目的地的关系。
    \item 法院受理案件后，应当适用哪国的实体法，\qd{只解决法律适用问题}。
    \item 标准结构。表述为\zd{什么问题适用什么法}。
          \begin{enumerate}
            \item 侵权（范围）/依/（关联词）侵权行为地（\zd{连结点}）/法（系属）。
            \item 连连结点都记不住剥离不出来找不着准据法。
            \item \zd{由 xxx 法院管辖×}。\qd{不要和管辖权规范混淆。永远管辖权在先法律适用在后}。
          \end{enumerate}
    \item 类型。\zd{判断很重要，判断错了做题就错}。
          \begin{enumerate}
            \item 单边。一个连结点且已明示国别。
            \item 双边（第二多）。一个连结点，结合案情才能确定国别。
            \item 重叠。两个以上连结点，且须同时适用（并列项）。\al{适用收养人和被收养人经居地法}。
            \item 选择（最多）。选择其中之一。无条件选择（A法 \hz\, B法）；有条件选择（有顺序要求，有一先一，没一再二）。
          \end{enumerate}
    \item 区际法律冲突（不同区域适用不同法律）。直接适用\zd{最密切联系原则}。\al{\scx{首先使用首都，再最密切}}。
  \end{enumerate}
\end{kk}
